% SPDX-License-Identifier: Apache-2.0
% covers: AxiPerPins, AxiPerDriven
\datasheet{AxiPerPins}{The AXI link's peripheral end joined to an AXI4 peripheral's pins.}

\dsfacts{\code{lib/parts/src/bus/axi\_per\_pins.rs}}%
{\code{txhdl\_parts::bus::axi\_per\_pins::AxiPerPins}}%
{\code{//docs:pcie}}

\subsection*{Function}

A core from elsewhere that is an AXI4 peripheral, such as AMD's DDR3
controller generated with its AXI4 port, has pins: a \code{valid} and a
\code{ready} per channel and a wire per field. \code{AxiPerPins} is
\code{AxiPins} the other way round. It stands on the link's peripheral
end, where an \code{AxiPer} would, and puts the link's channels onto
those pins. What it holds is what it owes the peripheral: the one write
it is giving, kept through the design's reset, which a peripheral such
as AMD's controller does not share, and counts of the responses and
read beats outstanding (issue 1532).

Its inputs are the side \code{AxiPerPinsIn}, destructured: the design's
reset, \code{rst}; the eleven pins the peripheral drives as one field,
\code{pins}, an \code{AxiPerDriven}; and the \code{aw}, \code{ar} and
\code{w} channels from the link. So its input ports are \code{rst},
\code{pins\_awready} and the rest, \code{aw}, \code{ar} and \code{w}.
\code{AxiPerDriven} derives \code{Ports}, so a design holding a
peripheral's pins among its own ports passes them on whole. Its outputs
are the side \code{outp: AxiPerPinsOut}: the \code{b} and \code{r}
channels back to the link, and the twenty-six pins the peripheral reads.

A reset can land after the peripheral took a write's address and
before it took the beats, which lived in the link's channel and went
with the reset. The peripheral would then pair every later write's
beats with the address before it, and a write's response would never
come. So the unit keeps the write it is giving the peripheral, one at
a time: \code{cur} says one is being given, \code{held} that its
address is not yet taken, \code{h\_id} to \code{h\_qos} hold that
address as the pins carry it, \code{beats} its beats and \code{sent}
how many have gone. \code{bowe} counts the writes taken whose response
has not come, \code{rowe} the read beats asked for that have not come,
and \code{drain} says the unit is draining after a reset.

\code{axi\_per\_pins::sim} has \code{PinRam}, a memory on the pins
written as a simulation, and \code{pins}, which makes the wires.
\code{PinRam} can take a warm-up, before which it is not ready, and a
read and a write latency, and it can take an address past its end
modulo its size. With \code{reading\_at\_address()} it takes a read's
words when it takes the read's address, rather than when it offers
each beat, as a controller that reads the memory before it answers
does (issue 1424). A write that lands between the two is then not in
the read's words. The link allows that, since it orders a read against
nothing written on the other channel. Those three are how
\code{ddr3::Ddr3} models AMD's controller.

\subsection*{Parameters}

\begin{center}\footnotesize
\begin{tabular}{@{}lp{0.7\textwidth}@{}}
\toprule
Parameter & Meaning \\
\midrule
\code{A} & The address width, in bits. The tables use 32. \\
\code{D} & The data width, in bits. The tables use 32. \\
\code{S} & The strobe width, \code{D / 8}. The tables use 4. \\
\code{I} & The identifier width, in bits. The tables use 5. \\
\bottomrule
\end{tabular}
\end{center}

The tables are generated for \code{AxiPerPins<32, 32, 4, 5>}, the
widths \code{Ddr3Per} uses it with on the board. The example
\code{ex\_axi\_per\_pins} uses \code{AxiPerPins<30, 32, 4, 5>}, the
controller's own 30-bit address.

\dstables{AxiPerPins}

\subsection*{Behaviour}

\begin{itemize}
\item A write starts when none is being given, an address phase is at
the head of \code{aw} and the unit is not draining. Its address phase
is \code{awvalid} and the \code{aw} pins: \code{awid}, \code{awaddr},
\code{awlen}, \code{awsize}, the burst type encoded in two bits,
\code{awlock}, \code{awcache}, \code{awprot} and \code{awqos}. It is
taken off the channel in the step it starts and kept in \code{h\_id}
to \code{h\_qos}, which drive the pins from then on, so the offer
stands until \code{awready} takes it. One write follows another with
no step between them (issue 1121).
\item The write's beats come from \code{w}, \code{wvalid} with
\code{wdata}, \code{wstrb} and \code{wlast}, until \code{sent} reaches
its count of beats, \code{awlen} plus one; a beat is taken off
\code{w} in a step \code{wready} is high. The write is done once its
address is taken and its beats have all gone.
\item An address phase at the head of \code{ar} is \code{arvalid} and
the \code{ar} pins, unless the unit is draining, and is taken off the
channel in a step \code{arready} is high.
\item A response the peripheral offers, \code{bvalid} with
\code{bid} and \code{bresp}, the response decoded from its two bits,
goes onto \code{b} in a step \code{b} has room, and \code{bready} is
that room. The same holds for \code{rvalid} and \code{r}, with
\code{rid}, \code{rdata}, \code{rresp} and \code{rlast}, and
\code{rready}.
\item \code{bowe} counts up as the peripheral takes a write's address
and down as a response is taken; \code{rowe} counts up by a read's
beats as the peripheral takes its address and down as a beat is taken.
A response or a read beat counts only while one is owed: at power-up a
controller's pins are unknown in a simulation of the netlist, and the
counts have no reset to wash that out.
\item Under \code{rst}, and after it until no write is being given and
nothing is owed, the unit drains. It goes on offering the kept
address, gives the rest of the write's beats with no byte enabled,
which writes nothing, with \code{wlast} on the last of its count, and
takes nothing from \code{w}. \code{bready} and \code{rready} are high,
so the responses and read beats meant for the design that was reset
are taken and dropped, none passed on to the link. It offers the
peripheral nothing new.
\item Nothing it offers waits on a \code{ready}, as AXI4 asks of a
host, and an offer is never withdrawn.
\item Every field of the link's address phase reaches a pin except
\code{region}: AMD's controller has no region pins, and a peripheral
joined to one link needs none, since the region is how an interconnect
tells decoded ranges of one peripheral apart. \code{AxQOS} is forwarded.
\end{itemize}

\subsection*{Verification}

\begin{itemize}
\item \code{an\_answer\_nothing\_asked\_for\_is\_not\_counted} in
\code{bus::axi\_per\_pins}, in \code{//lib/parts:parts\_test}, has the
pins offer a response and a read beat under the reset and after it,
with nothing asked: \code{bowe} and \code{rowe} stay zero and the drain
ends.
\item \code{a\_host\_on\_the\_link\_reaches\_a\_peripheral\_on\_pins}
there writes
three words through an \code{AxiHost}, the unit and a \code{PinRam},
reads four back across the burst's ends, and reads past the memory's
end, answered \code{SlvErr}.
\item \code{a\_timed\_memory\_waits\_and\_then\_answers\_late} there
holds \code{PinRam}'s timing: a warm-up holds a write off, and a latency
delays an answer by itself and no more.
\item \code{a\_memory\_reading\_at\_the\_address}\allowbreak
\code{\_lets\_a\_read\_pass\_a\_write}
there sends a read that waits twenty cycles and then a write to its
word that lands while it waits: the read gets the old word with
\code{reading\_at\_address()} and the new word without it.
\item \code{ex\_axi\_per\_pins} runs \code{AxiPerPins<30, 32, 4, 5>} with
a host on the link and a \code{PinRam} on the pins, and the build
simulates its netlist against that run under nvc and Verilator:
\code{//docs:axi\_per\_pins\_sim\_axi\_per\_pins\_tb\_test} and
\code{//docs:axi\_per\_pins\_vsim\_test}.
\item \code{//ddr3:ddr3\_test} drives it inside \code{Ddr3Per}, and
\code{//ddr3/sim:wrapper\_test}, which is manual, checks the
controller's side of the same pins under Vivado's simulator.
\item \code{the\_loader\_takes\_an\_image\_after\_a\_reset}\allowbreak
\code{\_in\_the\_middle\_of\_a\_program}, in
\code{//cpu/vreteno:board\_test}, resets the board's design from the
serial line while a program copies in the DDR3, sends on the Ethernet
port and has Razboj draw, at three distances into the program, and
asserts that one landed with a write's beats owed. The DDR3
controller, which that reset does not reach, keeps what it was doing,
and the loader must take the next image whole: acknowledge every
block, read back the sum it was sent, and run it.
\end{itemize}

\subsection*{Used by}

\code{Ddr3Per}, as its field \code{pins}, between the board's router
and AMD's DDR3 controller. The example \code{ex\_axi\_per\_pins}.

\subsection*{Limits}

Outside the drain, its \code{ready} outputs are the channels' room in
the same step, so a path runs from a channel's state through the unit
to the peripheral's pins with no register. It gives the peripheral one
write at a time: the next write's address is offered only once the
last one's address is taken and its beats have gone. It has no
\code{AxREGION} or user pins, so a peripheral that wants a region gets
it as a constant where it is instantiated.
